\section{Cohomological frame: $R_w$ as an Alexander polynomial}

We have produced three objects --- the cocycle $S$, the reduced polynomial $P_w$, and the core
$R_w$. This section fixes their cohomological meaning: $S$ is a $1$-cocycle on $F_2$, $P_w$ is
its restriction to the loop $\langle w\rangle$, and $R_w$ is the \emph{Alexander polynomial} of
the one-relator group $\langle u,v\mid w\rangle$ --- the order of an $H^2$.

\subsection*{The cocycle and its restriction}

By Proposition~\ref{prop:cocycle}, $S=\chi\circ\partial_u$ is a $1$-cocycle, so
$[S]\in H^1(F_2;\Lam)$ for the coefficients $\Lam$ with $F_2$ acting through $\chi$. As $F_2$ is
free it has cohomological dimension $1$, whence $H^n(F_2;-)=0$ for $n\ge2$: there is no $H^2$
over $F_2$ itself, and the relevant class lives in $H^1$.

Restrict to the cyclic subgroup $\langle w\rangle\cong\Z$. For $\Z$ acting on a module $M$
through $\tau:=\chi(w)$ one has $H^1(\Z;M)=M/(\tau-1)M$ (the coinvariants), a $1$-cocycle being
determined by its value on the generator modulo $(\tau-1)M$. Hence:

\begin{prop}\label{prop:restrict}
The restriction of $[S]$ to $\langle w\rangle$ is
\[
   \operatorname{res}_{\langle w\rangle}[S]\;=\;[\,S(w)\,]\;=\;P_w
   \ \in\ H^1(\langle w\rangle;\Lam)=\Lam/(1-\chi(w)),
\]
the class of $S(w)$ modulo the closure ideal (Proposition~\ref{prop:conj}). So $P_w$ is the
cocycle class evaluated on the loop $w$.
\end{prop}

\subsection*{The Alexander matrix of the relator}

Alongside $S=\chi\circ\partial_u$ put the companion cocycle $T:=\chi\circ\partial_v$; on positive
words $T(w)=\sum_{k:\,x_k=v}X^{\alpha_k}Y^{\beta_k}$ (Proposition~\ref{prop:closed} with $u,v$
exchanged).

\begin{lem}[Crowell--Fox fundamental formula]\label{lem:cf}
$\partial_u w\,(u-1)+\partial_v w\,(v-1)=w-1$ in $\Z[F_2]$; applying $\chi$,
\[
   (X-1)\,S(w)+(Y-1)\,T(w)=\chi(w)-1 ,
\]
so $(X-1)S(w)+(Y-1)T(w)\equiv0$ in $\Lam/(1-\chi(w))$. (Certified in \texttt{gen\_table.py}.)
\end{lem}

\subsection*{The one-relator group and $H^2$}

Let $G_w:=\langle u,v\mid w\rangle$. For coprime type the necklace is aperiodic, so $w$ is not a
proper power in $F_2$; by Lyndon's theorem the presentation $2$-complex is aspherical, a
$K(G_w,1)$. The character $\chi$ descends to $G_w$ exactly in $M:=\Lam/(1-\chi(w))$, where
$\chi(w)=1$, making $M$ a $G_w$-module (for coprime type the image of $\chi$ is the infinite
cyclic group $\langle t\rangle$, so $M=\Z[t,t^{-1}]=\Z[\Z]$). The cellular cochain complex of the
presentation complex (one $0$-cell, two $1$-cells $u,v$, one $2$-cell $w$),
\[
   M\;\xrightarrow{\ d^0\ }\;M^2\;\xrightarrow{\ d^1\ }\;M,\qquad
   d^1(m_u,m_v)=S(w)\,m_u+T(w)\,m_v,
\]
has top map the transpose of the Fox Jacobian $\partial_2=(\partial_u w,\partial_v w)$ --- the
\emph{Alexander matrix} of the relator. Hence
\[
   H^2(G_w;M)=\operatorname{coker}(d^1)=M/\big(S(w),T(w)\big).
\]

\subsection*{$R_w$ is the Alexander polynomial}

\begin{thm}\label{thm:alex}
For a coprime word $w$, over $\Q[t^{\pm1}]$,
\[
   \gcd\big(S(w),\,T(w)\big)\;=\;R_w .
\]
Consequently the elementary ideal $(S(w),T(w))\otimes\Q=(R_w)$, and
\[
   \Q\otimes H^2(G_w;M)\;\cong\;\Q[t^{\pm1}]/(R_w).
\]
Thus $R_w$ is the \emph{Alexander polynomial} $\Delta_{(G_w,\chi)}$ --- the order of the
(rational) Alexander module --- and $\deg R_w=\dim_\Q\big(\Q\otimes H^2(G_w;M)\big)$.
\end{thm}

\begin{proof}
Write $P_w=S(w)$, $T_w=T(w)$ in $M=\Z[t^{\pm1}]$, and $g_k:=1+t+\cdots+t^{k-1}=(t^k-1)/(t-1)$.
Clearing denominators in Lemma~\ref{lem:cf} (with $X=t^{-(E-o)}$, $Y=t^{o}$) gives
$(t^{E-o}-1)P_w=t^{E-o}(t^{o}-1)T_w$, i.e.\ $g_{E-o}\,P_w=t^{E-o}g_o\,T_w$. By
Corollary~\ref{cor:Rone}, $P_w=g_o R_w$; substituting and cancelling $g_o$ yields
$T_w=t^{-(E-o)}g_{E-o}R_w$ --- a unit times $g_{E-o}R_w$. Therefore
\[
   \gcd(P_w,T_w)=R_w\cdot\gcd(g_o,g_{E-o}).
\]
The common roots of $g_o$ and $g_{E-o}$ would be common nontrivial $o$-th and $(E-o)$-th roots
of unity; since $\gcd(o,E-o)=\gcd(E,o)=1$ there are none, so $\gcd(g_o,g_{E-o})=1$ and
$\gcd(P_w,T_w)=R_w$. Over the principal ideal domain $\Q[t^{\pm1}]$ the ideal $(P_w,T_w)$ is then
$(R_w)$, whence $\Q\otimes H^2(G_w;M)=\Q[t^{\pm1}]/(R_w)$.
\end{proof}

\subsection*{Dictionary}

\begin{rem}\label{rem:dict}
Theorem~\ref{thm:alex} reads the whole catalogue as Alexander-module data.
\begin{itemize}
\item[(i)] $\Mahler(R_w)$ is the Mahler measure of an Alexander polynomial. Lehmer's number is
the Mahler measure of the Alexander polynomial of the $(-2,3,7)$-pretzel knot; its appearance at
type $(12,7)$ (Proposition~\ref{prop:lehmer}) is the cocycle realisation of that classical fact.
\item[(ii)] The reflection symmetry of Proposition~\ref{prop:cores} is \emph{external}:
complement-then-reverse exchanges $\partial_u\leftrightarrow\partial_v$ (i.e.\ $S\leftrightarrow T$)
and $t\mapsto t^{-1}$, and as $\gcd$ is symmetric it carries $R_w$ to the same core. Individual
cores need \emph{not} be self-reciprocal --- unlike a knot's $\Delta$ --- see
Remark~\ref{rem:duality}.
\item[(iii)] $P_w=g_o R_w$ (Corollary~\ref{cor:Rone}) is the splitting ``(forced cyclotomic
part, killed in passing from $H^1(\langle w\rangle)$ to $H^2(G_w)$) $\times$ (Alexander core
$R_w$)''.
\item[(iv)] The four cycles of Section~2 carry their Alexander polynomials: $R_w=1$ for
$\{-1\},\{1\}$; $R_w=\Phi_2$ for $-5$; and $R_w=t^{11}-t^{10}+t^{7}-t^{6}+t^{4}-t+1$ for $-17$
(so $\dim_\Q H^2=11$).
\end{itemize}
\end{rem}

\subsection*{The Fox--Lyndon complex, torsion, and the deeper strata}

\begin{rem}[The resolution and its torsion]\label{rem:torsion}
The objects above are the low-degree pieces of one complex. The presentation $2$-complex of
$G_w$ being aspherical, its universal cover gives a length-$2$ free resolution of $\Z$ over
$\Z[G_w]$,
\[
   0\to \Z[G_w]\xrightarrow{\,(\partial_u w,\,\partial_v w)\,}\Z[G_w]^2
        \xrightarrow{\,(u-1,\,v-1)\,}\Z[G_w]\xrightarrow{\ \varepsilon\ }\Z\to0,
\]
with the Fox Jacobian as its only syzygy (Fox's free calculus; Lyndon). Tensoring to the abelian
cover $\Lam_w:=\Z[H_1(G_w)]$ (for coprime type $\Lam_w=\Z[t^{\pm1}]$) yields the \emph{Alexander
chain complex}
\[
   \Lam_w\xrightarrow{\,(S,\,T)\,}\Lam_w^2
         \xrightarrow{\,(t^{-(E-o)}-1,\ t^{o}-1)\,}\Lam_w,\qquad
   \partial_1\partial_2=0\ \text{(Lemma~\ref{lem:cf})},
\]
with $H_0=\Z$, $H_2=0$, and $H_1=A_w$ the Alexander module, of order $R_w$ --- the homological
mirror of $H^2=M/(S,T)$ in Theorem~\ref{thm:alex} (Crowell). Since $R_w\ne0$ the complex is
acyclic over $\Q(t)$, hence carries a Reidemeister--Milnor torsion; by Milnor's formula
\[
   \tau(G_w,\chi)\;\doteq\;\frac{R_w(t)}{t-1}\ \in\ \Q(t)^\times/\!\pm t^{\Z}
\]
(the $t-1$ being the order of $H_0=\Lam_w/(t-1)$). Thus $R_w$ is the order of $H_1$ and the
numerator of the torsion --- the master invariant of which the catalogue computes the polynomial
part. (Here $H_2=\ker(S,T)=0$, as $\Lam_w$ is a domain and $S,T\ne0$; and the ascending chain
of Fitting ideals of $A_w$, whose first term is the order ideal $(R_w)$ rationally, gives
strictly finer invariants than the polynomial alone.)
\end{rem}

\begin{rem}[The non-coprime stratum is torsion in $H_1$]\label{rem:noncoprime}
The first homology $H_1(G_w)=G_w^{\mathrm{ab}}=\Z^2/\langle(o,E-o)\rangle$ is read off the Smith
normal form of the relator vector $[\,o,\ E-o\,]$, namely $[\,d,\ 0\,]$ with $d=\gcd(E,o)$:
\[
   H_1(G_w)\;\cong\;\Z\oplus\Z/d .
\]
Coprime type is exactly $d=1$ --- the relator class $(o,E-o)$ primitive, $H_1$ free of rank one
on $t$. For $d>1$ the $\Z/d$ summand is the formal variable $s$ of \S1--\S2 with $s^{d}=1$, and
\[
   \Q\otimes\Z[H_1(G_w)]\;=\;\Q[t^{\pm1},s]/(s^{d}-1)\;\cong\;\prod_{e\mid d}\Q(\zeta_e)[t^{\pm1}],
\]
which is exactly the Chinese-remainder splitting of \S2. So the $s$-grading is the torsion of the
relator's homology class --- not a second free variable, since $H_1$ has rank one. Accordingly
$R_w$ is the $e=1$ (trivial-character) component of the $\Z/d$-twisted Alexander module over
$\Z[H_1(G_w)]=\Z[t^{\pm1},s]/(s^{d}-1)$, and specialising $s\mapsto\zeta_e$ --- the very programme
of \S1 --- is its isotypic decomposition over the $\Z/d$-cover, the components $e>1$ being the
genuinely new data of the non-coprime stratum.
\end{rem}

\begin{rem}[Reciprocity and the duality locus]\label{rem:duality}
For a knot the Alexander polynomial is symmetric, $\Delta(t)\doteq\Delta(t^{-1})$, forced by
Poincar\'e duality of the complement and refined by a Blanchfield linking pairing
$A_w\times A_w\to\Q(t)/\Lam_w$. \emph{This fails here in general}: the cores are usually not
self-reciprocal --- e.g.\ the $-17$ core $t^{11}-t^{10}+t^{7}-t^{6}+t^{4}-t+1$ is not --- because
the $2$-complex of $G_w$ is not a closed manifold. What survives is the \emph{external} symmetry
of Proposition~\ref{prop:cores}: complement-then-reverse sends $w$ to a word with the same core,
$R_{\overline{w^{R}}}=R_w$, via $S\leftrightarrow T$ and $t\mapsto t^{-1}$. The self-reciprocal
cores --- those fixed by reciprocation, $R_w(t)\doteq R_w(t^{-1})$ --- form a distinguished
\emph{duality locus}: of the $1529$ distinct cores with $E\le16$, exactly $107$ are
self-reciprocal --- the $55$ cyclotomic cores (all reciprocal) and $52$ non-cyclotomic ones. The
latter are the reciprocal cores, every \emph{Salem} core among them (Salem polynomials are
reciprocal), the smallest being Lehmer's polynomial; the non-reciprocal cores, like the $-17$
core, lie off the locus. A sufficient condition to land on it is that $w$ be a palindrome up to
rotation ($w^{R}\sim w$): then $R_w=R_{w^{R}}\doteq R_w(t^{-1})$. Whether the words on this locus
carry an honest Blanchfield-type sesquilinear pairing on $A_w$ --- a true duality --- is the
natural open question.
\end{rem}
